\subsubsection{Maze Navigation}
Figure~\ref{fig:qual_result}(a) shows agent visitation heatmaps in Maze Navigation across different difficulty levels.
In \emph{Easy} and \emph{Medium}, the density concentrates near the shortest paths, while retaining some pre-convergence spread, indicating a reliable goal-reaching policy.
Consistent with corridor geometry, $\textsc{Boost}$ is used early to accelerate along straight segments while staying within the per-episode budget.
In \emph{Hard}, the agent can encounter deadlocks, which can be resolved by activating the $\textsc{Penetration}$ option at chokepoints.
Fig.~\ref{fig:qual_result}(b) further illustrates such deadlock cases.
Overall, the heatmap patterns align with the quantitative metrics, showing that $\textsc{Boost}$ enables faster progress along straight corridors and $\textsc{Penetration}$ resolves blockages under budget constraints.

\subsubsection{Air-to-Air Combat}
TART exhibits both strong representation ability and adaptive combat behaviors. 
In Fig.~\ref{fig:main-figure}(a), maneuvers differ with or without missile launch, showing causal dependency between discrete and continuous actions.
Figure~\ref{fig:main-figure}(b) demonstrates multi-modality, where distinct VQ codes yield different maneuvers under the same conditions.
Figures~\ref{fig:qual_result}(c)-(d) illustrate adaptive maneuvers in dynamic combat scenarios.
In (c), the agent executes an offensive maneuver with consecutive $\textsc{Missile}$ actions.
In (d), it performs a defensive maneuver, neutralizing the opponent's missile with a $\textsc{Defense}$ action and following with a counter shot.
These results show that integrating TART enables the RL agent to capture causal dependencies and multi-modality while coordinating offensive and defensive options in a context-aware manner.